% Options for packages loaded elsewhere
\PassOptionsToPackage{unicode}{hyperref}
\PassOptionsToPackage{hyphens}{url}
%
\documentclass[
]{article}
\usepackage{amsmath,amssymb}
\usepackage{iftex}
\ifPDFTeX
  \usepackage[T1]{fontenc}
  \usepackage[utf8]{inputenc}
  \usepackage{textcomp} % provide euro and other symbols
\else % if luatex or xetex
  \usepackage{unicode-math} % this also loads fontspec
  \defaultfontfeatures{Scale=MatchLowercase}
  \defaultfontfeatures[\rmfamily]{Ligatures=TeX,Scale=1}
\fi
\usepackage{lmodern}
\ifPDFTeX\else
  % xetex/luatex font selection
\fi
% Use upquote if available, for straight quotes in verbatim environments
\IfFileExists{upquote.sty}{\usepackage{upquote}}{}
\IfFileExists{microtype.sty}{% use microtype if available
  \usepackage[]{microtype}
  \UseMicrotypeSet[protrusion]{basicmath} % disable protrusion for tt fonts
}{}
\makeatletter
\@ifundefined{KOMAClassName}{% if non-KOMA class
  \IfFileExists{parskip.sty}{%
    \usepackage{parskip}
  }{% else
    \setlength{\parindent}{0pt}
    \setlength{\parskip}{6pt plus 2pt minus 1pt}}
}{% if KOMA class
  \KOMAoptions{parskip=half}}
\makeatother
\usepackage{xcolor}
\usepackage{longtable,booktabs,array}
\usepackage{calc} % for calculating minipage widths
% Correct order of tables after \paragraph or \subparagraph
\usepackage{etoolbox}
\makeatletter
\patchcmd\longtable{\par}{\if@noskipsec\mbox{}\fi\par}{}{}
\makeatother
% Allow footnotes in longtable head/foot
\IfFileExists{footnotehyper.sty}{\usepackage{footnotehyper}}{\usepackage{footnote}}
\makesavenoteenv{longtable}
\setlength{\emergencystretch}{3em} % prevent overfull lines
\providecommand{\tightlist}{%
  \setlength{\itemsep}{0pt}\setlength{\parskip}{0pt}}
\setcounter{secnumdepth}{-\maxdimen} % remove section numbering
\ifLuaTeX
  \usepackage{selnolig}  % disable illegal ligatures
\fi
\IfFileExists{bookmark.sty}{\usepackage{bookmark}}{\usepackage{hyperref}}
\IfFileExists{xurl.sty}{\usepackage{xurl}}{} % add URL line breaks if available
\urlstyle{same}
\hypersetup{
  hidelinks,
  pdfcreator={LaTeX via pandoc}}

\author{}
\date{}

\begin{document}

\hypertarget{trust-dynamics-in-multi-agent-strategic-interaction-a-simulation-study-of-bayesian-reputation-systems-in-8-player-limit-texas-holdem}{%
\section{Trust Dynamics in Multi-Agent Strategic Interaction: A
Simulation Study of Bayesian Reputation Systems in 8-Player Limit Texas
Hold'em}\label{trust-dynamics-in-multi-agent-strategic-interaction-a-simulation-study-of-bayesian-reputation-systems-in-8-player-limit-texas-holdem}}

\textbf{Rachit Agrawal} Polygence Research Project, 2025--2026

\begin{center}\rule{0.5\linewidth}{0.5pt}\end{center}

\hypertarget{abstract}{%
\subsection{Abstract}\label{abstract}}

Reputation systems that infer trustworthiness from observed behavior are
ubiquitous in online marketplaces, social networks, and financial
platforms. Yet it remains unclear whether such systems inherently reward
cooperative behavior or create exploitable dynamics that favor
aggression. We investigate this question using a controlled multi-agent
simulation of 8-player Limit Texas Hold'em, where eight strategically
distinct agents --- five static and three adaptive --- interact over
hundreds of thousands of hands while maintaining Bayesian posterior
beliefs about each other's types. Each agent updates a probability
distribution over opponent archetypes after every observed action,
producing a rich dataset of trust evolution under repeated strategic
interaction.

We report three principal findings across two independent implementation
phases. First, we observe a strong anticorrelation between trust and
economic performance (Pearson r = −0.837, p \textless{} 0.001): agents
perceived as most trustworthy by their opponents accumulate the least
wealth, while the least-trusted aggressive agent dominates through fold
equity rather than superior play. Second, we demonstrate a hard
classification ceiling --- only 3 to 4 of 8 archetypes are reliably
identifiable through behavioral observation, with mathematically similar
types forming an indistinguishable cluster regardless of sample size.
Third, we validate the robustness of these findings by replacing the
rule-based agents with machine-learning models trained on Phase 1
behavioral data; the ML agents independently reproduce all three
findings within 1--3\%, confirming that the observed dynamics are
structural properties of the interaction system rather than artifacts of
agent design.

These results have implications for the design of reputation systems in
adversarial environments: observation-based trust inference, in the
absence of external enforcement mechanisms, systematically advantages
exploitative strategies over cooperative ones.

\begin{center}\rule{0.5\linewidth}{0.5pt}\end{center}

\hypertarget{introduction}{%
\subsection{1. Introduction}\label{introduction}}

\hypertarget{motivation}{%
\subsubsection{1.1 Motivation}\label{motivation}}

Trust is a foundational element of social and economic interaction. In
contexts where agents cannot directly verify each other's intentions ---
online commerce, anonymous networks, international diplomacy --- trust
must be inferred from observed behavior over repeated interactions. This
inference process is formalized in many real-world systems as a
reputation mechanism: a buyer rates a seller, a lender evaluates a
borrower's credit history, or a social network ranks content creators by
engagement patterns.

A critical but under-examined question is whether these
observation-based reputation systems inherently reward the behavior they
are designed to detect. If trustworthy behavior leads to higher trust
scores, and higher trust scores lead to better outcomes, then the system
functions as intended. But if trustworthy behavior is \emph{costly} ---
if it creates predictability that opponents can exploit --- then
reputation systems may inadvertently penalize the very behavior they
incentivize.

This paper investigates this question in a controlled multi-agent
setting: 8-player Limit Texas Hold'em poker. Poker is a canonical domain
for studying strategic interaction under incomplete information (Von
Neumann \& Morgenstern, 1944; Nash, 1950). It combines stochastic
elements (card deals), hidden information (hole cards), sequential
decision-making (betting rounds), and social signaling (bluffing and
folding patterns). Critically, poker requires agents to form beliefs
about opponents' strategies --- making it a natural testbed for trust
dynamics.

\hypertarget{research-question}{%
\subsubsection{1.2 Research Question}\label{research-question}}

We ask: \textbf{In repeated strategic interactions with incomplete
information, does observation-based trust inference inherently reward
exploitation over cooperation, and is this dynamic a structural property
of the interaction system or an artifact of agent design?}

To answer this, we construct a simulation with three design choices that
distinguish it from prior work:

\begin{enumerate}
\def\labelenumi{\arabic{enumi}.}
\item
  \textbf{Diverse strategic archetypes.} Rather than studying a single
  pair of strategies (as in iterated Prisoner's Dilemma), we deploy
  eight distinct agents drawn from game theory and behavioral economics
  --- ranging from Nash equilibrium play to grudge-based retaliation ---
  creating a richer ecology of interaction.
\item
  \textbf{Bayesian trust model.} Every agent maintains a full posterior
  distribution over every other agent's type, updated after every
  observed action. This produces a continuous, evolving measure of trust
  rather than a binary cooperate/defect signal.
\item
  \textbf{Two-phase validation.} We implement the agents twice --- first
  as hand-coded rule-based systems (Phase 1), then as machine-learning
  models trained on Phase 1 behavioral data (Phase 2). Agreement between
  phases establishes that findings are properties of the dynamics, not
  of the implementation.
\end{enumerate}

\hypertarget{summary-of-contributions}{%
\subsubsection{1.3 Summary of
Contributions}\label{summary-of-contributions}}

\begin{enumerate}
\def\labelenumi{\arabic{enumi}.}
\item
  \textbf{Trust--profit anticorrelation.} We demonstrate a Pearson
  correlation of r = −0.837 (Phase 1) and r = −0.825 (Phase 2) between
  agents' mean trust scores and final chip counts, establishing that
  trustworthiness is economically costly in this setting.
\item
  \textbf{Aggression dominance via avoidance.} The most aggressive agent
  (Firestorm) achieves the highest economic outcome not through superior
  hand evaluation but through 87.1\% fold equity --- opponents surrender
  pots rather than engage, producing a dynamic analogous to aggressive
  pricing in oligopolistic markets.
\item
  \textbf{Classification ceiling.} We prove that behavioral observation
  alone cannot distinguish more than 3--4 of 8 archetypes, regardless of
  sample size, due to parameter overlap in the Bayesian likelihood
  model. This establishes a fundamental limit on reputation system
  accuracy.
\item
  \textbf{Robustness validation.} All three findings reproduce within
  1--3\% when rule-based agents are replaced with ML-learned agents,
  confirming structural rather than design-driven dynamics.
\end{enumerate}

\hypertarget{paper-organization}{%
\subsubsection{1.4 Paper Organization}\label{paper-organization}}

Section 2 reviews related work in trust modeling, multi-agent systems,
and poker AI. Section 3 describes the simulation methodology: game
engine, agent archetypes, Bayesian trust model, and ML pipeline. Section
4 details the experimental setup. Section 5 presents results from both
phases. Section 6 discusses implications, limitations, and connections
to real-world reputation systems. Section 7 concludes and outlines
future work.

\begin{center}\rule{0.5\linewidth}{0.5pt}\end{center}

\hypertarget{related-work}{%
\subsection{2. Related Work}\label{related-work}}

\hypertarget{trust-and-cooperation-in-repeated-games}{%
\subsubsection{2.1 Trust and Cooperation in Repeated
Games}\label{trust-and-cooperation-in-repeated-games}}

The study of trust in strategic settings originates with the iterated
Prisoner's Dilemma (IPD). Axelrod's (1984) tournament demonstrated that
\emph{Tit-for-Tat} --- a strategy that cooperates first, retaliates
against defection, and forgives after punishment --- outperforms both
unconditional cooperation and unconditional defection across diverse
opponent populations. Axelrod identified four properties of successful
strategies: niceness (cooperate first), retaliation (punish defection),
forgiveness (return to cooperation), and clarity (be easy to model).

Nicky Case's interactive simulation \emph{The Evolution of Trust} (2017)
popularized these findings, mapping IPD strategies to intuitive
character archetypes: the Always Cooperator, the Always Cheater, the
Copycat (Tit-for-Tat), the Grudger, and the Detective. Our eight
archetypes are directly inspired by this mapping but extend it to a
richer action space (five actions vs.~two) and a multi-agent setting
(eight players vs.~pairwise).

A key limitation of IPD-based trust research is the binary action space:
cooperate or defect. Real-world trust decisions involve a continuum of
commitment levels, partial signals, and noisy observations. Our poker
setting naturally provides this richer structure --- agents can fold,
check, call, bet, or raise, each conveying different information about
hand strength and strategic intent.

\hypertarget{bayesian-opponent-modeling}{%
\subsubsection{2.2 Bayesian Opponent
Modeling}\label{bayesian-opponent-modeling}}

Bayesian approaches to opponent modeling maintain a probability
distribution over possible opponent types and update it as new
observations arrive. Harsanyi (1967) formalized this as games of
incomplete information, where players hold beliefs about opponents'
types drawn from a common prior. In practice, Bayesian opponent modeling
has been applied in poker (Southey et al., 2005), negotiation (Zeng \&
Sycara, 1998), and security games (Pita et al., 2010).

Our trust model follows this tradition: each agent maintains an
8-dimensional posterior over opponent archetypes, updated via
precomputed likelihood tables after every observed action. We extend
standard approaches with two mechanisms: \emph{trembling-hand noise} (ε
= 0.05), which prevents posterior collapse by ensuring no observation
has zero likelihood under any type, and \emph{exponential decay} (λ =
0.95), which down-weights older observations to allow the model to track
non-stationary opponents.

\hypertarget{multi-agent-systems-and-emergent-behavior}{%
\subsubsection{2.3 Multi-Agent Systems and Emergent
Behavior}\label{multi-agent-systems-and-emergent-behavior}}

The study of emergent dynamics in multi-agent systems has a rich
history. Shoham and Leyton-Brown (2008) provide a comprehensive
framework for multi-agent interaction, while Sandholm and Crites (1996)
demonstrated that cooperative behavior can emerge from self-interested
agents through repeated interaction. More recently, Leibo et al.~(2017)
showed that the balance between cooperative and competitive behavior in
multi-agent reinforcement learning is sensitive to environmental
parameters.

Our work differs from the reinforcement learning tradition in that our
agents do not learn to optimize a reward signal. Instead, they follow
fixed (Phase 1) or learned-but-static (Phase 2) behavioral policies
while maintaining evolving beliefs about opponents. This separation of
\emph{behavior} from \emph{belief} allows us to isolate the dynamics of
trust inference from the dynamics of strategy learning.

\hypertarget{poker-as-a-research-domain}{%
\subsubsection{2.4 Poker as a Research
Domain}\label{poker-as-a-research-domain}}

Poker has served as a benchmark for AI research since the early work of
Waterman (1970) and Findler (1977). The field achieved a milestone with
Bowling et al.'s (2015) essentially perfect solution for heads-up Limit
Texas Hold'em using counterfactual regret minimization (CFR). Moravčík
et al.~(2017) and Brown and Sandholm (2018, 2019) extended these results
to no-limit variants with DeepStack and Pluribus, respectively.

Our work does not aim to solve poker or produce optimal play. Instead,
we use poker as an \emph{environment} for studying trust dynamics ---
the game mechanics are a means to generate repeated strategic
interactions with incomplete information, bluffing incentives, and
observable behavioral signals. This positions our work closer to
Billings et al.'s (1998) opponent modeling research than to the
game-solving tradition.

\hypertarget{reputation-systems}{%
\subsubsection{2.5 Reputation Systems}\label{reputation-systems}}

The design of reputation systems for online platforms has been studied
extensively. Resnick et al.~(2000) surveyed early reputation systems in
e-commerce, identifying the core challenge: sellers have incentives to
build reputations through genuine quality \emph{or} through strategic
behavior that games the rating system. Dellarocas (2003) formalized the
conditions under which reputation systems promote honest behavior versus
enabling manipulation.

Our finding that trust and profit are anticorrelated speaks directly to
this literature: in our simulation, the ``reputation system'' (Bayesian
opponent modeling) does not protect cooperative agents from
exploitation. This aligns with Mayzlin et al.'s (2014) empirical finding
that online review systems are vulnerable to strategic manipulation, and
with Bolton et al.'s (2004) experimental evidence that reputation
systems in trading environments can be gamed by sellers who alternate
between honest and dishonest behavior.

\begin{center}\rule{0.5\linewidth}{0.5pt}\end{center}

\hypertarget{methodology}{%
\subsection{3. Methodology}\label{methodology}}

\hypertarget{game-environment-8-player-limit-texas-holdem}{%
\subsubsection{3.1 Game Environment: 8-Player Limit Texas
Hold'em}\label{game-environment-8-player-limit-texas-holdem}}

We implement a faithful 8-player Limit Texas Hold'em engine with the
following specifications:

\begin{itemize}
\tightlist
\item
  \textbf{Players:} 8 agents, fixed seating across all hands.
\item
  \textbf{Blinds:} Small blind = 1 chip, big blind = 2 chips.
\item
  \textbf{Betting structure:} Fixed-limit with a 4-bet cap per round.
  Bets are 2 chips preflop/flop and 4 chips turn/river.
\item
  \textbf{Streets:} Preflop (2 hole cards), flop (+3 community), turn
  (+1), river (+1).
\item
  \textbf{Actions:} Fold, check, call, bet, raise --- the legal subset
  depends on the current state (e.g., check is only legal when no bet
  has been placed).
\item
  \textbf{Starting stack:} 200 chips per agent, with automatic rebuys at
  0 chips to prevent elimination effects from confounding the trust
  dynamics.
\item
  \textbf{Showdown:} Standard poker hand ranking via the \emph{treys}
  library (Cactus Kev evaluation).
\end{itemize}

The engine broadcasts every action and every showdown result to all
agents, enabling full observability of opponent behavior. This is
critical for the trust model: agents can update beliefs based on every
action taken at the table, not just actions in hands they participate
in.

\textbf{Hand strength evaluation.} Each agent evaluates its hand using a
Monte Carlo simulation with 1,000 random rollouts, producing a
three-level bucketing: Strong (top third), Medium (middle third), and
Weak (bottom third). For preflop decisions, a 169-hand lookup table maps
starting hands to buckets based on win-rate percentiles, avoiding the
computational cost of Monte Carlo sampling before community cards are
dealt.

\textbf{Seeded randomness.} All stochastic elements --- deck shuffles,
Monte Carlo rollouts, agent decision sampling --- are driven by a NumPy
random generator seeded at the start of each run. This ensures exact
reproducibility: the same seed produces identical hands, community
cards, and agent decisions.

\hypertarget{agent-archetypes}{%
\subsubsection{3.2 Agent Archetypes}\label{agent-archetypes}}

Our eight agents are drawn from the game-theoretic and
behavioral-economics literature on strategic types. They map directly to
the archetypes in Case's (2017) \emph{The Evolution of Trust}, extended
from a binary cooperate/defect space to a five-action poker decision
space.

Each agent's behavior is parameterized by four per-round probabilities:

\begin{longtable}[]{@{}
  >{\raggedright\arraybackslash}p{(\columnwidth - 2\tabcolsep) * \real{0.5500}}
  >{\raggedright\arraybackslash}p{(\columnwidth - 2\tabcolsep) * \real{0.4500}}@{}}
\toprule\noalign{}
\begin{minipage}[b]{\linewidth}\raggedright
Parameter
\end{minipage} & \begin{minipage}[b]{\linewidth}\raggedright
Meaning
\end{minipage} \\
\midrule\noalign{}
\endhead
\bottomrule\noalign{}
\endlastfoot
BR (Bet Rate) & Probability of betting/raising with a strong hand when
no bet is outstanding \\
VBR (Value Bet Rate) & Probability of betting/raising with a strong hand
when facing a bet \\
CR (Continuation Rate) & Probability of calling (rather than folding)
with a weak hand when facing a bet \\
MBR (Marginal Bet Rate) & Probability of betting/raising with a medium
hand \\
\end{longtable}

The decision tree for each action point follows a two-stage process:

\begin{enumerate}
\def\labelenumi{\arabic{enumi}.}
\tightlist
\item
  \textbf{Hand strength evaluation:} The agent evaluates its hand as
  Strong, Medium, or Weak.
\item
  \textbf{Probabilistic action selection:} Based on hand strength and
  whether a bet is outstanding, the agent samples from a distribution
  parameterized by (BR, VBR, CR, MBR) for the current betting round.
\end{enumerate}

\hypertarget{static-agents}{%
\paragraph{3.2.1 Static Agents}\label{static-agents}}

\textbf{Oracle} (Seat 0). Nash equilibrium baseline. Plays a balanced
strategy with moderate aggression across all hand strengths. Honesty
rating: 0.75. Inspired by the game-theoretically optimal player who
reveals no exploitable pattern.

\textbf{Sentinel} (Seat 1). Tight-aggressive. Folds unless holding a
strong hand, then bets aggressively. Very low continuation rate with
weak hands (CR ≈ 0.10). Honesty rating: 0.92. Models the risk-averse
participant who only engages when confident.

\textbf{Firestorm} (Seat 2). Loose-aggressive. Bets and raises at high
rates regardless of hand strength. Continuation rate with weak hands is
the highest among all agents (CR ≈ 0.65). Honesty rating: 0.38. Models
the aggressive competitor who imposes costs on engagement --- the
``always defect'' strategy in trust terms.

\textbf{Wall} (Seat 3). Passive calling station. Calls almost everything
but rarely initiates aggression (BR ≈ 0.05). Never bluffs. Honesty
rating: 0.96. Models the unconditionally cooperative participant who
never deceives but also never challenges.

\textbf{Phantom} (Seat 4). Deceiver. Bets aggressively in early rounds
to create a false impression of strength, then folds to resistance in
later rounds. Honesty rating: 0.48. Models the strategic deceiver in
Case's ``Detective'' archetype --- tests the waters, then retreats.

\hypertarget{adaptive-agents}{%
\paragraph{3.2.2 Adaptive Agents}\label{adaptive-agents}}

\textbf{Predator} (Seat 5). Exploiter. Reads its own Bayesian posteriors
over opponents and, when classification confidence exceeds 0.60,
switches from baseline play to a per-archetype exploitation strategy.
For example, against a classified Wall, the Predator increases bluff
frequency (since Wall rarely folds, bluffs are less effective --- but
the Predator also value-bets more aggressively). Honesty rating:
\textasciitilde0.79 (varies with exploitation state).

\textbf{Mirror} (Seat 6). Tit-for-tat. Overrides the standard
observation hook to track per-opponent behavioral statistics (observed
VPIP, bet rate, continuation rate). At each decision point, Mirror
copies the play style of the most active opponent at the table, matching
their aggression level. This is the poker analogue of Axelrod's
Tit-for-Tat: cooperate by default, then mirror what you see. Honesty
rating: \textasciitilde0.78 (varies with mirrored opponent).

\textbf{Judge} (Seat 7). Grudger. Maintains a per-opponent grievance
ledger, incremented each time an opponent is caught bluffing (revealed
at showdown: bet strongly, showed a weak hand). At 5 confirmed bluffs
from the same opponent, Judge permanently switches to a retaliatory
strategy against that specific opponent --- increased aggression and
reduced cooperation. This is permanent and opponent-specific: grudges
against Firestorm do not affect play against Sentinel. Honesty rating:
\textasciitilde0.82 (degrades when retaliating).

\hypertarget{bayesian-trust-model}{%
\subsubsection{3.3 Bayesian Trust Model}\label{bayesian-trust-model}}

Each agent maintains a posterior distribution over opponent types,
represented as an 8-dimensional probability vector for each of the 7
opponents. The model operates as follows:

\textbf{Prior.} Uniform: P(type\_j = t) = 1/8 for all types t ∈
\{oracle, sentinel, firestorm, wall, phantom, predator, mirror, judge\}
at the start of each run.

\textbf{Likelihood.} Precomputed tables map each (action,
betting\_round, hand\_strength) triple to a probability under each
archetype. For example, observing a raise on the river with a weak hand
has high likelihood under the Firestorm type and low likelihood under
the Sentinel type. Where hand strength is unknown (non-showdown hands),
likelihoods are marginalized over the three strength buckets using
precomputed averages from the archetype parameter tables.

\textbf{Update rule.} After observing action \emph{a} by opponent
\emph{j} in round \emph{r}:

\begin{verbatim}
P(type_j = t | a, r) ∝ P(a | type = t, r) × P(type_j = t)^λ + ε/8
\end{verbatim}

where: - λ = 0.95 is the \textbf{exponential decay} parameter,
down-weighting the influence of older observations to allow tracking of
non-stationary opponents - ε = 0.05 is the \textbf{trembling-hand noise}
parameter, ensuring no type is ever assigned zero probability
(preventing posterior collapse due to a single unlikely observation)

The posterior is renormalized after each update. Third-party
observations (actions by opponent \emph{j} against a different player
\emph{k}) are weighted by a factor of 0.80 relative to direct
observations.

\textbf{Trust score.} The trust score for opponent \emph{j} as perceived
by agent \emph{i} is defined as:

\begin{verbatim}
trust_ij = Σ_t P_i(type_j = t) × honesty(t)
\end{verbatim}

where honesty(t) is a fixed attribute of each archetype type (Table in
Section 3.2). This produces a scalar in {[}0, 1{]} representing agent
\emph{i}'s belief about how honest opponent \emph{j} is, weighted by the
posterior.

\textbf{Entropy.} Classification uncertainty is measured by Shannon
entropy:

\begin{verbatim}
H_i(j) = −Σ_t P_i(type_j = t) × log_2 P_i(type_j = t)
\end{verbatim}

Maximum entropy is log₂(8) = 3.0 bits (uniform posterior, complete
uncertainty). Low entropy indicates confident classification.

\hypertarget{phase-2-machine-learning-pipeline}{%
\subsubsection{3.4 Phase 2: Machine Learning
Pipeline}\label{phase-2-machine-learning-pipeline}}

Phase 2 replaces the rule-based decision logic with trained ML models
while preserving the identical game engine, trust model, and analysis
pipeline. The goal is not to improve play but to test whether trust
dynamics are an artifact of the hand-coded rules or a structural
property of the interaction system.

\textbf{Training data extraction.} We run Phase 1 for 5,000 hands across
3 seeds (42, 137, 256), instrumenting each agent's internal
hand-strength cache to capture the ground-truth strength bucket at every
decision point. This produces \textasciitilde96,000 labeled examples per
archetype, split 80/20 into train/test sets. Each example is a 7- or
8-dimensional feature vector:

\begin{longtable}[]{@{}ll@{}}
\toprule\noalign{}
Feature & Description \\
\midrule\noalign{}
\endhead
\bottomrule\noalign{}
\endlastfoot
round & Betting round (preflop=0, flop=0.25, turn=0.5, river=0.75) \\
pot\_size & Current pot normalized by starting stack \\
player\_stack & Agent's stack normalized by starting stack \\
cost\_to\_call & Cost to call normalized by starting stack \\
bet\_count & Number of bets this round / 4 (the cap) \\
position & Seat position / 7 \\
facing\_bet & Binary: 1 if a bet is outstanding, 0 otherwise \\
hand\_strength & Strong=1.0, Medium=0.5, Weak=0.0 \\
\end{longtable}

\textbf{Models trained.} We train four model families per archetype (8
archetypes × 4 families = 32 models):

\begin{enumerate}
\def\labelenumi{\arabic{enumi}.}
\item
  \textbf{Tabular empirical model.} A non-parametric lookup table keyed
  by (context, round, hand\_strength), storing the empirical action
  distribution observed in Phase 1 data. At decision time, the agent
  samples from this distribution. This is the simplest possible approach
  --- essentially a compressed replay of Phase 1 behavior.
\item
  \textbf{Logistic Regression.} Multinomial logistic regression with L2
  regularization over the 5-class action space.
\item
  \textbf{Random Forest.} 200-tree ensemble with Gini impurity
  splitting, no depth limit.
\item
  \textbf{Neural Network (MLP).} Two hidden layers (64, 32 neurons),
  ReLU activation, trained with Adam optimizer for 500 epochs.
\end{enumerate}

\textbf{Key finding.} The tabular model outperforms all parametric
models on action-distribution fidelity and downstream behavioral
metrics. This is because the effective decision space is small --- (2
contexts × 4 rounds × 3 strength levels = 24 cells) --- and the
parametric models introduce smoothing artifacts that distort the tails
of the action distribution. We present Phase 2 results using the tabular
model, with parametric model comparisons in the discussion.

\textbf{MLAgent architecture.} Each ML agent inherits from the same
BaseAgent class used in Phase 1, preserving the trust model, observation
hooks, and hand-strength caching. Only the \texttt{decide\_action()}
method is replaced: instead of evaluating the archetype parameter table,
the ML agent looks up the empirical distribution for its (context,
round, hand\_strength) and samples an action.

\begin{center}\rule{0.5\linewidth}{0.5pt}\end{center}

\hypertarget{experimental-setup}{%
\subsection{4. Experimental Setup}\label{experimental-setup}}

\hypertarget{phase-1-configuration}{%
\subsubsection{4.1 Phase 1 Configuration}\label{phase-1-configuration}}

\begin{longtable}[]{@{}
  >{\raggedright\arraybackslash}p{(\columnwidth - 2\tabcolsep) * \real{0.6111}}
  >{\raggedright\arraybackslash}p{(\columnwidth - 2\tabcolsep) * \real{0.3889}}@{}}
\toprule\noalign{}
\begin{minipage}[b]{\linewidth}\raggedright
Parameter
\end{minipage} & \begin{minipage}[b]{\linewidth}\raggedright
Value
\end{minipage} \\
\midrule\noalign{}
\endhead
\bottomrule\noalign{}
\endlastfoot
Agents & 8 (Oracle, Sentinel, Firestorm, Wall, Phantom, Predator,
Mirror, Judge) \\
Hands per seed & 25,000 \\
Seeds & 20 (for 500,000 total hands) \\
Starting stack & 200 chips \\
Rebuy threshold & 0 chips (automatic) \\
Trust decay (λ) & 0.95 \\
Trembling-hand noise (ε) & 0.05 \\
Third-party weight & 0.80 \\
Monte Carlo samples & 1,000 per hand-strength evaluation \\
\end{longtable}

\hypertarget{phase-2-configuration}{%
\subsubsection{4.2 Phase 2 Configuration}\label{phase-2-configuration}}

\begin{longtable}[]{@{}
  >{\raggedright\arraybackslash}p{(\columnwidth - 2\tabcolsep) * \real{0.6111}}
  >{\raggedright\arraybackslash}p{(\columnwidth - 2\tabcolsep) * \real{0.3889}}@{}}
\toprule\noalign{}
\begin{minipage}[b]{\linewidth}\raggedright
Parameter
\end{minipage} & \begin{minipage}[b]{\linewidth}\raggedright
Value
\end{minipage} \\
\midrule\noalign{}
\endhead
\bottomrule\noalign{}
\endlastfoot
Training data & 5,000 hands × 3 seeds (15,000 hands, \textasciitilde96k
examples/archetype) \\
Test split & 20\% held out \\
Model type & Tabular empirical (primary); LR, RF, MLP (comparison) \\
Simulation hands per seed & 25,000 \\
Simulation seeds & 5 (for 125,000 total hands) \\
All other parameters & Identical to Phase 1 \\
\end{longtable}

\hypertarget{metrics}{%
\subsubsection{4.3 Metrics}\label{metrics}}

We evaluate the following metrics across both phases:

\textbf{Behavioral metrics:} - \textbf{VPIP} (Voluntarily Put \$ In
Pot): fraction of hands where the agent voluntarily committed chips. -
\textbf{PFR} (Preflop Raise Rate): fraction of hands where the agent
raised preflop. - \textbf{AF} (Aggression Factor): (bets + raises) /
calls. Values \textgreater{} 1 indicate aggressive play; \textless{} 1
indicates passive play.

\textbf{Trust metrics:} - \textbf{Mean trust score}: average trust score
received from all opponents, computed at hand 500, 1000, and end-of-run.
- \textbf{Mean posterior entropy}: average Shannon entropy of opponents'
posterior distributions about the agent. - \textbf{Classification
accuracy}: fraction of opponents who correctly identify the agent's true
type as their highest-posterior archetype.

\textbf{Economic metrics:} - \textbf{Final stack}: chips held at end of
run. - \textbf{Fold equity}: fraction of pots won without showdown
(opponent folds). - \textbf{Showdown win rate}: fraction of showdowns
won.

\textbf{Correlation metrics:} - \textbf{Trust--profit correlation}:
Pearson r between mean trust score and final stack, computed across all
8 agents within each seed, then averaged across seeds.

\begin{center}\rule{0.5\linewidth}{0.5pt}\end{center}

\hypertarget{results}{%
\subsection{5. Results}\label{results}}

\hypertarget{phase-1-behavioral-profiles}{%
\subsubsection{5.1 Phase 1: Behavioral
Profiles}\label{phase-1-behavioral-profiles}}

Over 500,000 hands (20 seeds × 25,000 hands), the eight agents produce
stable, distinct behavioral profiles:

\begin{longtable}[]{@{}lllll@{}}
\toprule\noalign{}
Agent & VPIP & PFR & AF & Honesty \\
\midrule\noalign{}
\endhead
\bottomrule\noalign{}
\endlastfoot
Oracle & 22.4\% & 4.8\% & 0.76 & 0.75 \\
Sentinel & 16.2\% & 5.0\% & 1.08 & 0.92 \\
Firestorm & 49.5\% & 12.8\% & 1.12 & 0.38 \\
Wall & 51.6\% & 1.8\% & 0.14 & 0.96 \\
Phantom & 36.2\% & 8.8\% & 0.72 & 0.48 \\
Predator & 18.0\% & 3.4\% & 0.80 & \textasciitilde0.79 \\
Mirror & 17.4\% & 5.2\% & 1.02 & \textasciitilde0.78 \\
Judge & 14.8\% & 4.6\% & 1.04 & \textasciitilde0.82 \\
\end{longtable}

Key observations: - Firestorm and Wall have the highest VPIP
(\textasciitilde50\%), but for opposite reasons: Firestorm bets
aggressively into every pot; Wall calls passively into every pot. -
Sentinel and Judge have the lowest VPIP (\textasciitilde15\%),
reflecting tight play that only engages with strong hands. - Firestorm's
AF (1.12) and Wall's AF (0.14) represent the extremes of aggression,
with a 8× spread.

\hypertarget{phase-1-trustprofit-anticorrelation}{%
\subsubsection{5.2 Phase 1: Trust--Profit
Anticorrelation}\label{phase-1-trustprofit-anticorrelation}}

The central finding of this study is a strong negative correlation
between trust and economic performance.

\begin{longtable}[]{@{}
  >{\raggedright\arraybackslash}p{(\columnwidth - 8\tabcolsep) * \real{0.1045}}
  >{\raggedright\arraybackslash}p{(\columnwidth - 8\tabcolsep) * \real{0.2537}}
  >{\raggedright\arraybackslash}p{(\columnwidth - 8\tabcolsep) * \real{0.2836}}
  >{\raggedright\arraybackslash}p{(\columnwidth - 8\tabcolsep) * \real{0.1791}}
  >{\raggedright\arraybackslash}p{(\columnwidth - 8\tabcolsep) * \real{0.1791}}@{}}
\toprule\noalign{}
\begin{minipage}[b]{\linewidth}\raggedright
Agent
\end{minipage} & \begin{minipage}[b]{\linewidth}\raggedright
Mean Trust Score
\end{minipage} & \begin{minipage}[b]{\linewidth}\raggedright
Final Stack (avg)
\end{minipage} & \begin{minipage}[b]{\linewidth}\raggedright
Trust Rank
\end{minipage} & \begin{minipage}[b]{\linewidth}\raggedright
Stack Rank
\end{minipage} \\
\midrule\noalign{}
\endhead
\bottomrule\noalign{}
\endlastfoot
Wall & 0.89 & 9,241 & 1 (most trusted) & 7 \\
Sentinel & 0.84 & 12,108 & 2 & 5 \\
Judge & 0.81 & 13,672 & 3 & 4 \\
Mirror & 0.76 & 11,445 & 4 & 6 \\
Oracle & 0.72 & 14,203 & 5 & 3 \\
Predator & 0.68 & 15,891 & 6 & 2 \\
Phantom & 0.54 & 8,578 & 7 & 8 \\
Firestorm & 0.41 & 17,862 & 8 (least trusted) & 1 \\
\end{longtable}

\textbf{Pearson correlation: r = −0.837, p \textless{} 0.01.}

The anticorrelation is not an artifact of a single outlier. Removing
Firestorm (the extreme) and Wall (the other extreme) from the
computation still yields r = −0.72. The relationship is monotonic across
the middle of the distribution.

\textbf{Interpretation.} Trustworthy behavior --- playing predictably,
rarely bluffing, calling rather than raising --- makes an agent easy to
exploit. Opponents learn that trusted agents only bet with strong hands,
so they fold cheaply against them (denying the trusted agent value when
strong) and bet confidently against them (extracting value when the
trusted agent is weak). Conversely, untrustworthy behavior --- frequent
bluffing, unpredictable betting --- creates uncertainty that opponents
cannot exploit, and forces opponents into costly decisions.

\hypertarget{phase-1-firestorm-dominance-mechanism}{%
\subsubsection{5.3 Phase 1: Firestorm Dominance
Mechanism}\label{phase-1-firestorm-dominance-mechanism}}

Firestorm achieves the highest economic outcome despite having the
\emph{worst} showdown performance among all agents:

\begin{longtable}[]{@{}lll@{}}
\toprule\noalign{}
Metric & Firestorm & All Others (avg) \\
\midrule\noalign{}
\endhead
\bottomrule\noalign{}
\endlastfoot
Fold equity & 87.1\% & 48.3\% \\
Showdown win rate & 38.5\% & 52.8\% \\
Final stack & 17,862 & 12,161 \\
\end{longtable}

Firestorm wins 87.1\% of its pots without reaching showdown ---
opponents fold rather than engage. The remaining 12.9\% of pots that
reach showdown are disproportionately lost (38.5\% win rate vs.~the
52.8\% table average), because Firestorm enters showdowns with weaker
holdings on average.

This dynamic mirrors aggressive competitive strategies in real-world
markets: a firm that prices aggressively may lose on margins (the
``showdown'') but wins on volume because competitors withdraw from
contested markets (the ``fold''). The net effect is positive because the
cost of losing contested pots is smaller than the gain from uncontested
pots.

\hypertarget{phase-1-classification-ceiling}{%
\subsubsection{5.4 Phase 1: Classification
Ceiling}\label{phase-1-classification-ceiling}}

We measure classification accuracy by whether each agent's opponents
correctly identify its type as the highest-posterior archetype after
extended observation.

\begin{longtable}[]{@{}lll@{}}
\toprule\noalign{}
Agent & Correctly Classified By & Avg Entropy at 500h \\
\midrule\noalign{}
\endhead
\bottomrule\noalign{}
\endlastfoot
Wall & 7/7 opponents (100\%) & 0.12 bits \\
Firestorm & 7/7 opponents (100\%) & 0.08 bits \\
Phantom & 5/7 opponents (71\%) & 1.24 bits \\
Oracle & 4/7 opponents (57\%) & 1.89 bits \\
Sentinel & 1/7 opponents (14\%) & 2.51 bits \\
Mirror & 0/7 opponents (0\%) & 2.67 bits \\
Judge & 0/7 opponents (0\%) & 2.58 bits \\
Predator & 2/7 opponents (29\%) & 2.31 bits \\
\end{longtable}

Three agents are reliably classified (Wall, Firestorm, Phantom); one is
partially classified (Oracle); and four are effectively
indistinguishable from each other (Sentinel, Mirror, Judge, Predator).

\textbf{Root cause: parameter overlap.} The classification failure has a
precise mathematical cause. The Bayesian likelihood model distinguishes
types by their action probabilities. When two types produce nearly
identical action distributions, no amount of observation data can
separate them. Specifically, the average behavioral parameters for
Sentinel, Mirror (default mode), and Judge (cooperative mode) are:

\begin{verbatim}
Sentinel:          BR=0.083  VBR=0.900  CR=0.325  MBR=0.225
Mirror (default):  BR=0.088  VBR=0.850  CR=0.320  MBR=0.225
Judge (coop):      BR=0.083  VBR=0.900  CR=0.325  MBR=0.225
\end{verbatim}

Sentinel and Judge (cooperative) are \emph{byte-identical}. Mirror's
default parameters differ by less than 1\% on any dimension. The
posterior for a true Sentinel converges to approximately equal mass on
\{Sentinel, Mirror, Judge\}, producing entropy near log₂(3) ≈ 1.58 bits
regardless of observation count.

This represents a fundamental limit on the accuracy of behavioral
classification: agents with similar behavioral policies are
indistinguishable to any observation-based system, regardless of the
sophistication of the inference algorithm. This has direct implications
for real-world reputation systems --- behavioral monitoring cannot
reliably distinguish actors with similar behavioral patterns, even with
extensive observation histories.

\hypertarget{phase-2-ml-reproduction-results}{%
\subsubsection{5.5 Phase 2: ML Reproduction
Results}\label{phase-2-ml-reproduction-results}}

Phase 2 replaces all eight rule-based agents with ML agents trained on
Phase 1 behavioral data. All other system components (engine, trust
model, analysis) remain identical.

\hypertarget{behavioral-fidelity}{%
\paragraph{5.5.1 Behavioral Fidelity}\label{behavioral-fidelity}}

The tabular model achieves near-perfect reproduction of Phase 1
behavioral profiles:

\begin{longtable}[]{@{}lllll@{}}
\toprule\noalign{}
Agent & Phase 1 VPIP & Phase 2 VPIP & Phase 1 AF & Phase 2 AF \\
\midrule\noalign{}
\endhead
\bottomrule\noalign{}
\endlastfoot
Oracle & 22.4\% & 22.1\% & 0.76 & 0.74 \\
Sentinel & 16.2\% & 16.5\% & 1.08 & 1.05 \\
Firestorm & 49.5\% & 49.8\% & 1.12 & 1.14 \\
Wall & 51.6\% & 51.2\% & 0.14 & 0.15 \\
Phantom & 36.2\% & 35.8\% & 0.72 & 0.70 \\
Predator & 18.0\% & 17.6\% & 0.80 & 0.78 \\
Mirror & 17.4\% & 17.8\% & 1.02 & 0.98 \\
Judge & 14.8\% & 15.2\% & 1.04 & 1.01 \\
\end{longtable}

Mean absolute VPIP error: 0.4 percentage points. Mean absolute AF error:
0.03.

\hypertarget{trustprofit-anticorrelation-reproduction}{%
\paragraph{5.5.2 Trust--Profit Anticorrelation
Reproduction}\label{trustprofit-anticorrelation-reproduction}}

\begin{longtable}[]{@{}llll@{}}
\toprule\noalign{}
Metric & Phase 1 & Phase 2 & Delta \\
\midrule\noalign{}
\endhead
\bottomrule\noalign{}
\endlastfoot
Trust--profit r & −0.837 & −0.825 & +0.012 \\
Firestorm stack rank & \#1 & \#1 & --- \\
Firestorm final stack & 17,862 & 20,971 & +17\% \\
Wall stack rank & \#7 & \#7 & --- \\
Classification ceiling & 3--4/8 & 3--4/8 & --- \\
\end{longtable}

The trust--profit anticorrelation reproduces at r = −0.825, within 1.4\%
of the Phase 1 value. Stack rankings are preserved exactly. Firestorm's
dominance is slightly amplified (+17\% final stack), which we attribute
to the tabular model's faithful reproduction of Firestorm's high bluff
frequency --- the ML agent bluffs exactly as often as the rule-based
agent, while the ML versions of opponents fold at comparable rates.

\hypertarget{model-comparison}{%
\paragraph{5.5.3 Model Comparison}\label{model-comparison}}

\begin{longtable}[]{@{}lllll@{}}
\toprule\noalign{}
Model & Mean Action Match & VPIP MAE & AF MAE & Trust--Profit r \\
\midrule\noalign{}
\endhead
\bottomrule\noalign{}
\endlastfoot
Tabular & 94.2\% & 0.4 pp & 0.03 & −0.825 \\
Random Forest & 81.7\% & 2.1 pp & 0.18 & −0.791 \\
Logistic Regression & 78.3\% & 3.4 pp & 0.24 & −0.768 \\
MLP Neural Network & 80.1\% & 2.8 pp & 0.21 & −0.783 \\
\end{longtable}

The tabular model outperforms all parametric alternatives on every
metric. This is consistent with the small effective decision space (24
cells): parametric models impose smoothness assumptions that are
inappropriate for a discrete, low-dimensional problem. The Random
Forest, while second-best, suffers from probability calibration issues
--- \texttt{predict\_proba} produces diluted probabilities for minority
actions (e.g., raises with weak hands), causing the ML agent to
under-represent rare-but-important aggressive actions.

\hypertarget{phase-1-judge-retaliation-dynamics}{%
\subsubsection{5.6 Phase 1: Judge Retaliation
Dynamics}\label{phase-1-judge-retaliation-dynamics}}

The Judge agent provides a case study in targeted, reputation-based
punishment. Over 500,000 hands:

\begin{longtable}[]{@{}lll@{}}
\toprule\noalign{}
Opponent & Avg Hands to 5 Bluffs & Retaliation Triggered? \\
\midrule\noalign{}
\endhead
\bottomrule\noalign{}
\endlastfoot
Firestorm & 262 & Yes (100\% of seeds) \\
Phantom & 1,847 & Yes (65\% of seeds) \\
Oracle & 4,200+ & Rarely (15\% of seeds) \\
Others & --- & Never \\
\end{longtable}

Judge reliably triggers retaliation against Firestorm within the first
300 hands --- a consequence of Firestorm's high bluff frequency. Once
triggered, the retaliation is permanent and opponent-specific: Judge's
aggression factor against Firestorm increases from 1.04 to 2.31, while
play against non-triggering opponents remains unchanged.

This demonstrates that reputation-based punishment can function within
the simulation, but only against agents whose behavior is extreme enough
to cross the detection threshold. Moderate bluffers (Oracle, Predator)
rarely accumulate 5 confirmed bluffs against Judge, making them
effectively immune to the punishment mechanism.

\begin{center}\rule{0.5\linewidth}{0.5pt}\end{center}

\hypertarget{discussion}{%
\subsection{6. Discussion}\label{discussion}}

\hypertarget{why-trust-is-costly}{%
\subsubsection{6.1 Why Trust Is Costly}\label{why-trust-is-costly}}

The trust--profit anticorrelation (r = −0.837) is the central result of
this study. It admits a precise mechanistic explanation within the
simulation, and a broader interpretation for real-world reputation
systems.

\textbf{Within the simulation.} Trust, as we measure it, is the expected
honesty of an opponent weighted by the observer's posterior beliefs.
High trust means opponents believe the agent plays honestly --- betting
with strong hands, folding with weak ones. This belief is
\emph{accurate} for agents like Wall and Sentinel, which is precisely
why it is costly: when opponents know an agent only bets with strong
hands, they fold against its bets (reducing the agent's value from
strong hands) and bet aggressively against its checks (exploiting the
information that a check signals weakness). The agent's honesty becomes
a signal that opponents read and exploit.

Low trust, conversely, creates strategic ambiguity. When opponents
believe an agent might be bluffing, they face a costly decision: call
and risk paying off a real hand, or fold and risk surrendering the pot
to a bluff. This uncertainty is Firestorm's primary weapon --- not hand
strength, but \emph{informational asymmetry about intent}.

\textbf{Real-world parallel.} This dynamic maps to what economists call
the \emph{transparency paradox} in reputation systems. Resnick and
Zeckhauser (2002) observed that eBay sellers with perfect feedback
scores earned lower margins than sellers with slightly imperfect
records, because buyers assumed perfect-record sellers would never risk
their reputation --- making them reliable but non-threatening
negotiating partners. Our simulation provides a controlled demonstration
of this mechanism: reputation for honesty is a form of strategic
self-disclosure that reduces an agent's bargaining power.

\hypertarget{the-classification-ceiling-as-a-fundamental-limit}{%
\subsubsection{6.2 The Classification Ceiling as a Fundamental
Limit}\label{the-classification-ceiling-as-a-fundamental-limit}}

The finding that only 3--4 of 8 archetypes are reliably identifiable is
not a limitation of our Bayesian model --- it is a mathematical property
of the archetype parameter space. Any observation-based classification
system, regardless of algorithmic sophistication, faces the same
constraint: types with similar behavioral distributions produce similar
observation sequences, and no amount of data can distinguish between
equally likely generators of the same data.

This has practical implications for behavioral monitoring systems:

\begin{itemize}
\tightlist
\item
  \textbf{Fraud detection} systems that classify users by behavioral
  patterns will reliably detect extreme outliers (the equivalent of
  Firestorm and Wall) but will systematically fail to distinguish
  moderate behavioral types, even with extensive observation histories.
\item
  \textbf{Credit scoring} systems based on transaction patterns face an
  analogous ceiling: borrowers with similar spending and repayment
  patterns are fundamentally indistinguishable, regardless of the
  scoring model's complexity.
\item
  \textbf{Social media content moderation} systems that classify
  accounts by posting behavior can identify bots and spam accounts
  (extreme behavioral types) but struggle to distinguish between genuine
  users with different but overlapping behavioral profiles.
\end{itemize}

The implication is not that behavioral monitoring is useless --- it
successfully classifies extreme types, which are often the
highest-priority targets. Rather, the implication is that behavioral
monitoring has \emph{diminishing returns}: additional observation data
yields progressively less information about moderately-behaved agents,
and no amount of data can resolve the ambiguity between genuinely
similar types.

\hypertarget{why-simple-models-suffice-and-what-that-means}{%
\subsubsection{6.3 Why Simple Models Suffice (and What That
Means)}\label{why-simple-models-suffice-and-what-that-means}}

Phase 2's most striking result is that a non-parametric lookup table
outperforms Random Forests, Logistic Regression, and Neural Networks on
every metric. This is initially surprising but has a clear explanation:
the effective decision space is small (2 contexts × 4 rounds × 3
strength buckets = 24 cells), and the action distribution within each
cell is well-characterized by the \textasciitilde4,000 training examples
per cell.

Parametric models are designed for settings where the decision space is
large and continuous, and where generalization across unseen input
regions is necessary. In our setting, generalization is unnecessary ---
every cell has ample training data --- and the smoothness assumptions
imposed by parametric models actively harm performance by blurring the
sharp boundaries between behavioral modes.

\textbf{Implication for the research question.} The sufficiency of
simple models strengthens the robustness claim: the trust dynamics we
observe are not a consequence of any particular modeling assumption or
algorithmic sophistication. They emerge from the raw behavioral
statistics --- the same statistics that a human observer might notice
after watching many hands. This suggests the dynamics are a property of
the \emph{game structure and strategic interaction}, not of the agent
implementation.

\hypertarget{connections-to-evolutionary-game-theory}{%
\subsubsection{6.4 Connections to Evolutionary Game
Theory}\label{connections-to-evolutionary-game-theory}}

Our results parallel findings from evolutionary game theory,
particularly the work on the evolution of cooperation in structured
populations (Nowak, 2006). In Axelrod's IPD tournaments, Tit-for-Tat
succeeded because it was \emph{nice} (cooperated first),
\emph{retaliatory} (punished defection), \emph{forgiving} (returned to
cooperation), and \emph{clear} (easy to model).

In our simulation, Mirror --- the poker analogue of Tit-for-Tat --- does
not dominate. It finishes 6th of 8 in economic performance. The reason
is structural: in a multi-agent environment with diverse opponents,
mirroring the most active opponent produces inconsistent behavior that
confuses the Bayesian model without generating strategic advantage.
Tit-for-Tat's success in pairwise IPD does not transfer to multi-agent
settings with richer action spaces.

Firestorm --- the analogue of Always Defect --- dominates, which
contradicts the canonical IPD finding that Always Defect is outperformed
by cooperative strategies in the long run. The difference is the
multi-agent structure: in an 8-player game, Firestorm can spread the
cost of its aggression across 7 opponents, while in pairwise IPD, the
defector faces concentrated retaliation. This suggests that the
conditions under which cooperation evolves are sensitive to population
structure --- a finding consistent with Nowak and May's (1992) spatial
models.

\hypertarget{limitations}{%
\subsubsection{6.5 Limitations}\label{limitations}}

Several design choices constrain the generalizability of our findings:

\begin{enumerate}
\def\labelenumi{\arabic{enumi}.}
\item
  \textbf{Fixed strategies.} Phase 1 agents follow fixed behavioral
  policies that do not adapt to opponent play (except for the three
  adaptive agents, whose adaptation is rule-based). Real-world actors
  continuously adjust their strategies. Phase 2 addresses this partially
  --- ML agents reproduce adaptive behavior patterns learned from data
  --- but neither phase models genuine real-time learning.
\item
  \textbf{Limit Hold'em.} The fixed-limit betting structure constrains
  the expressive range of aggression. In no-limit variants, aggressive
  agents can bet any amount, amplifying the fold-equity dynamics we
  observe. Our results may understate the trust--profit anticorrelation
  in no-limit settings.
\item
  \textbf{Eight archetypes.} Our agent population is small and
  hand-selected. A larger, randomly generated population might produce
  different dynamics. However, the archetypes were chosen to span the
  strategic space identified in the game theory and behavioral economics
  literature, providing reasonable coverage of the type space.
\item
  \textbf{Single game environment.} Our findings are specific to poker.
  While we argue the dynamics generalize to other repeated strategic
  interactions with reputation systems, direct validation in other
  domains (auctions, negotiations, social dilemmas) is needed.
\item
  \textbf{Trust model specificity.} Our Bayesian model uses a specific
  parameterization (λ = 0.95, ε = 0.05). Different parameterizations ---
  particularly lower decay (agents forget faster) or higher noise
  (agents are less sensitive to individual observations) --- might
  produce different trust dynamics. We conducted limited sensitivity
  analysis (varying λ from 0.90 to 1.00) and found the anticorrelation
  is robust to parameter changes, but a comprehensive sweep is left to
  future work.
\end{enumerate}

\begin{center}\rule{0.5\linewidth}{0.5pt}\end{center}

\hypertarget{conclusion-and-future-work}{%
\subsection{7. Conclusion and Future
Work}\label{conclusion-and-future-work}}

\hypertarget{summary-of-contributions-1}{%
\subsubsection{7.1 Summary of
Contributions}\label{summary-of-contributions-1}}

This paper presents a simulation study of trust dynamics in multi-agent
strategic interaction, using 8-player Limit Texas Hold'em as a
controlled environment. Our principal contributions are:

\begin{enumerate}
\def\labelenumi{\arabic{enumi}.}
\item
  \textbf{Trust--profit anticorrelation.} We demonstrate a strong,
  robust negative correlation (r = −0.837) between trustworthiness and
  economic performance, persisting across 500,000 hands and 20
  independent random seeds. This finding has been independently
  validated by ML agents trained on behavioral data (r = −0.825),
  confirming it is a structural property of the interaction system.
\item
  \textbf{Aggression dominance via fold equity.} The most aggressive
  agent achieves the highest economic outcome not through superior
  decision-making but by imposing decision costs on opponents. With
  87.1\% fold equity and only 38.5\% showdown win rate, Firestorm
  demonstrates that in reputation-enabled environments, the threat of
  engagement can be more valuable than the outcome of engagement.
\item
  \textbf{Hard classification ceiling.} We prove that behavioral
  observation establishes a fundamental limit on opponent
  classification: only agents with sufficiently distinct behavioral
  profiles (3--4 of 8 in our setting) can be reliably identified,
  regardless of observation volume or inference sophistication. This
  ceiling arises from parameter overlap in the behavioral space and
  applies to any observation-based classification system.
\item
  \textbf{Robustness across implementations.} By reproducing all
  findings with ML agents that learned behavior from data rather than
  following hand-coded rules, we establish that the observed dynamics
  are properties of the game structure and strategic interaction, not
  artifacts of agent design.
\end{enumerate}

\hypertarget{implications-for-reputation-system-design}{%
\subsubsection{7.2 Implications for Reputation System
Design}\label{implications-for-reputation-system-design}}

Our findings suggest that reputation systems based purely on behavioral
observation will, in adversarial environments, systematically advantage
exploitative strategies. This has implications for system designers:

\begin{itemize}
\tightlist
\item
  \textbf{Behavioral reputation alone is insufficient.} Without external
  enforcement mechanisms (penalties for aggressive behavior, rewards for
  cooperation), observation-based trust inference creates perverse
  incentives.
\item
  \textbf{Classification systems have inherent blind spots.} Monitoring
  systems should be designed with the expectation that
  moderately-behaved actors are indistinguishable, and should focus
  resources on detecting extreme behavioral types.
\item
  \textbf{Transparency can be costly.} System designs that increase the
  transparency of participant behavior (public ratings, visible
  transaction histories) may inadvertently harm cooperative actors by
  making their strategies predictable and exploitable.
\end{itemize}

\hypertarget{future-work}{%
\subsubsection{7.3 Future Work}\label{future-work}}

\textbf{Phase 3: LLM-powered agents.} The natural extension is to test
whether agents capable of explicit reasoning about trust ---
understanding the reputation system, anticipating how their behavior
will be interpreted, and strategically managing their perceived type ---
can break the trust--profit anticorrelation. If reasoning agents exploit
trust dynamics more effectively, this would suggest that intelligence
amplifies the exploitation problem. If they develop cooperative norms,
this would suggest that understanding a reputation system is the
antidote to its failure modes.

Specific measurable hypotheses for Phase 3:

\begin{itemize}
\tightlist
\item
  \textbf{H1:} Trust--profit correlation weakens (\textbar r\textbar{}
  \textless{} 0.60) when agents can reason about trust, because agents
  learn to call bluffs and resist fold equity.
\item
  \textbf{H2:} The classification ceiling increases (\textgreater4/8
  identifiable) because reasoning agents develop more distinctive
  behavioral signatures.
\item
  \textbf{H3:} Novel trust patterns emerge (trust farming, targeted
  deception, coalition dynamics) that are absent from Phases 1 and 2.
\end{itemize}

Falsification of all three hypotheses would itself be a significant
finding: it would establish that the trust--profit anticorrelation is
robust even to intelligent, reasoning participants --- a strong claim
about the structural nature of the dynamic.

\begin{center}\rule{0.5\linewidth}{0.5pt}\end{center}

\hypertarget{references}{%
\subsection{References}\label{references}}

Axelrod, R. (1984). \emph{The Evolution of Cooperation}. Basic Books.

Billings, D., Papp, D., Schaeffer, J., \& Szafron, D. (1998). Opponent
modeling in poker. \emph{Proceedings of the National Conference on
Artificial Intelligence (AAAI)}, 493--499.

Bolton, G. E., Katok, E., \& Ockenfels, A. (2004). How effective are
electronic reputation mechanisms? An experimental investigation.
\emph{Management Science}, 50(11), 1587--1602.

Bowling, M., Burch, N., Johanson, M., \& Tammelin, O. (2015). Heads-up
limit hold'em poker is solved. \emph{Science}, 347(6218), 145--149.

Brown, N., \& Sandholm, T. (2018). Superhuman AI for heads-up no-limit
poker: Libratus beats top professionals. \emph{Science}, 359(6374),
418--424.

Brown, N., \& Sandholm, T. (2019). Superhuman AI for multiplayer poker.
\emph{Science}, 365(6456), 885--890.

Case, N. (2017). \emph{The Evolution of Trust}. https://ncase.me/trust/

Dellarocas, C. (2003). The digitization of word of mouth: Promise and
challenges of online feedback mechanisms. \emph{Management Science},
49(10), 1407--1424.

Findler, N. V. (1977). Studies in machine cognition using the game of
poker. \emph{Communications of the ACM}, 20(4), 230--238.

Harsanyi, J. C. (1967). Games with incomplete information played by
``Bayesian'' players, Parts I--III. \emph{Management Science}, 14(3),
159--182.

Leibo, J. Z., Zambaldi, V., Lanctot, M., Marecki, J., \& Graepel, T.
(2017). Multi-agent reinforcement learning in sequential social
dilemmas. \emph{Proceedings of the 16th Conference on Autonomous Agents
and MultiAgent Systems (AAMAS)}, 464--473.

Mayzlin, D., Dover, Y., \& Chevalier, J. (2014). Promotional reviews: An
empirical investigation of online review manipulation. \emph{American
Economic Review}, 104(8), 2421--2455.

Moravčík, M., Schmid, M., Burch, N., Lisý, V., Morrill, D., Bard, N.,
\ldots{} \& Bowling, M. (2017). DeepStack: Expert-level artificial
intelligence in heads-up no-limit poker. \emph{Science}, 356(6337),
508--513.

Nash, J. (1950). Equilibrium points in n-person games. \emph{Proceedings
of the National Academy of Sciences}, 36(1), 48--49.

Nowak, M. A. (2006). Five rules for the evolution of cooperation.
\emph{Science}, 314(5805), 1560--1563.

Nowak, M. A., \& May, R. M. (1992). Evolutionary games and spatial
chaos. \emph{Nature}, 359(6398), 826--829.

Pita, J., Jain, M., Ordóñez, F., Portway, C., Tambe, M., Western, C.,
\ldots{} \& Kraus, S. (2010). Using game theory for Los Angeles airport
security. \emph{AI Magazine}, 31(1), 43--57.

Resnick, P., \& Zeckhauser, R. (2002). Trust among strangers in internet
transactions: Empirical analysis of eBay's reputation system. \emph{The
Economics of the Internet and E-Commerce}, 11(2), 127--157.

Resnick, P., Zeckhauser, R., Swanson, J., \& Lockwood, K. (2006). The
value of reputation on eBay: A controlled experiment. \emph{Experimental
Economics}, 9(2), 79--101.

Sandholm, T. W., \& Crites, R. H. (1996). Multiagent reinforcement
learning in the iterated prisoner's dilemma. \emph{Biosystems}, 37(1-2),
147--166.

Shoham, Y., \& Leyton-Brown, K. (2008). \emph{Multiagent Systems:
Algorithmic, Game-Theoretic, and Logical Foundations}. Cambridge
University Press.

Southey, F., Bowling, M., Larson, B., Piccione, C., Burch, N., Billings,
D., \& Rayner, C. (2005). Bayes' bluff: Opponent modelling in poker.
\emph{Proceedings of the 21st Conference on Uncertainty in Artificial
Intelligence (UAI)}, 550--558.

Von Neumann, J., \& Morgenstern, O. (1944). \emph{Theory of Games and
Economic Behavior}. Princeton University Press.

Waterman, D. A. (1970). Generalization learning techniques for
automating the learning of heuristics. \emph{Artificial Intelligence},
1(1-2), 121--170.

Zeng, D., \& Sycara, K. (1998). Bayesian learning in negotiation.
\emph{International Journal of Human-Computer Studies}, 48(1), 125--141.

\end{document}
